\subsection{\revdeldec{ Financial Time-Series }Imputation\revadddec{: Problem Formulation and Characteristics}}
\label{sec:imputation}

\revdelblock{ \subsection{Problem Formulation and Characteristics}



}Financial time-series imputation is fundamentally formulated as a bidirectional, closed-format generation task. Given a multivariate sequence $\mathbf{X} \in \mathbb{R}^{T \times d}$ and a corresponding binary indicator mask $\mathbf{M} \in \{0, 1\}^{T \times d}$ where $m_{t,i} = 1$ denotes a missing value, the objective is to learn a conditional distribution $p(\mathbf{X}_{miss} | \mathbf{X}_{obs}; \theta)$ to accurately reconstruct the missing elements $\mathbf{X}_{miss} = \mathbf{X} \odot \mathbf{M}$ conditioned strictly on the observed elements $\mathbf{X}_{obs} = \mathbf{X} \odot (\mathbf{1} - \mathbf{M})$.

In financial contexts, the missingness mechanisms—\revdeldec{ Missing Completely At Random }\revadddec{missing completely at random }(MCAR), \revdeldec{ Missing At Random }\revadddec{missing at random }(MAR), and \revdeldec{ Missing Not At Random }\revadddec{missing not at random }(MNAR)\revadddec{, formalized by Rubin~\citep{rubin1976inference}}—present unique challenges. MNAR is particularly prevalent due to structural market events such as sudden trading halts during market crashes or asynchronous trading hours across global exchanges\revadddec{, a setting related to non-synchronous observation of multivariate financial processes~\citep{hayashi2005covariance}}. Consequently, imputation models must conquer two primary domain-specific hurdles\revdeldec{ : (1) managing }\revadddec{. First, models must manage }\textit{contiguous missingness}, where extended data outages corrupt long temporal windows\revdeldec{ , and (2) maintaining }\revadddec{. Second, models must maintain }strict \textit{cross-variable consistency} without introducing look-ahead bias, which is critical to ensure that historical backtests of trading strategies remain structurally valid.

\subsection{\revadddec{Imputation: }Statistical, Interpolation, and \revdeldec{ Machine Learning }\revadddec{Machine-Learning }Baselines}

Early methodologies treated financial imputation primarily as a \revdeldec{ matrix completion }\revadddec{matrix-completion~\citep{candes2009exact} }or interpolation problem rooted in \revdeldec{ strict econometric }\revadddec{explicit statistical }assumptions. The classic Chow-Lin method \citep{chow1971} remains a foundational statistical approach for temporal disaggregation, frequently integrated with LASSO regression\revadddec{~\citep{tibshirani1996regression} }to estimate high-frequency (e.g., monthly) values from low-frequency (e.g., quarterly) observations using auxiliary market indicators. Later, \revdeldec{ machine learning }\revadddec{machine-learning }techniques such as \revdeldec{ KNN, MICE}\revadddec{$k$-nearest-neighbour imputation~\citep{troyanskaya2001missing}, multiple imputation by chained equations (MICE)~\citep{vanbuuren1999multiple}}, and Random Forests\revdeldec{ (e.g., missForest }\revadddec{~\citep{breiman2001random}, including missForest~}\citep{stekhoven2012}\revdeldec{ ) }\revadddec{, }offered flexible, non-parametric alternatives for tabular data imputation. The heterogeneous-data adversarial learning network proposed for FinTech applications by Khuwaja et al. \citep{NDDP} is included here as a related data-driven baseline. Operating in a continuous functional space, Kernel \revdeldec{ Ridge Regression }\revadddec{ridge regression~\citebox{\citep{saunders1998dual,aronszajn1950theory}}, in a fixed-income application~}\citep{smith2020}\revadddec{, }estimates missing discount curves by projecting noisy bond prices into a \revdeldec{ Reproducing Kernel Hilbert Space }\revadddec{reproducing kernel Hilbert space }(RKHS). This framework minimizes a weighted least-squares objective regularized with exponential weighting to deliberately suppress wild extrapolations during periods of high volatility.

As datasets expanded, traditional machine learning models such as \revdeldec{ K-Nearest Neighbors (KNN), Multiple Imputation by Chained Equations (MICE)}\revadddec{$k$-nearest neighbours~\citep{troyanskaya2001missing}, MICE~\citep{vanbuuren1999multiple}}, and Random Forest\revdeldec{ imputation 
}\revadddec{/missForest imputation~\citebox{\citep{breiman2001random,stekhoven2012}} }were widely adopted. Random Forest imputation, in particular, leverages ensemble decision trees to dynamically predict missing values by capturing complex, non-linear feature interactions. \revdeldec{ Empirical studies demonstrate that in specific financial contexts---such as corporate distress prediction---Random Forest frameworks can surprisingly outperformadvanced deep learning alternatives, particularly when data sparsity is extreme (e.g., missing rates of 40--50\%)}\revadddec{Evidence from corporate-distress prediction indicates that tree-based ensembles can remain competitive with, and in some sparse-data settings outperform, deep learners~\citep{sideras2024bankruptcy}; the precise magnitude of the advantage is dataset-dependent}.

\subsection{\revadddec{Imputation: }Tensor Factorization and Deep Sequential Approaches}

To accommodate higher-dimensional non-linearities and multivariate temporal dependencies, research progressed toward deep interpolation and tensor decomposition networks. Deep sequential fusion architectures, such as the BiGRU-BiLSTM framework proposed by \revadddec{Wang et al.~}\citep{wang2021bigru}, capture bidirectional temporal flows \revadddec{using gated recurrent units~\citep{cho2014learning} within a bidirectional recurrent architecture~\citep{schuster1997bidirectional}, }and employ a \revdeldec{ Group Method of Data Handling }\revadddec{group method of data handling }(GMDH) polynomial \revdeldec{ regression decoder}\revadddec{decoder~\citep{ivakhnenko1971polynomial} }to reconstruct complex missing patterns in multivariate series.

Concurrently, tensor factorization models emerged as highly efficient solutions for structural missingness. NMTucker~\citep{varolgunes2023nmtucker}\revadddec{, building on Tucker decomposition~\citep{tucker1966some}, }introduced a non-linear Matryoshka Tucker decomposition, systematically replacing linear core operations with neural network-based reconstructions and enforcing recursive core constraints to mitigate overfitting on sparse financial tensors. Expanding upon this, RegTensor \citep{zhou2023regtensor} facilitates fast, non-linear coupled tensor completion across multiple sparse datasets. By sharing factor matrices across tensors, it captures latent cross-market dynamics while strictly enforcing uncorrelated temporal embeddings via an orthogonal regularization penalty:
\begin{equation}
\mathcal{L}_{RegTensor} = \mathcal{L}_{recon} + \lambda \|\mathbf{U}^\top \mathbf{U} - \mathbf{I}\|_F^2
\end{equation}
where $\mathbf{U}$ represents the temporal factor matrix and $\lambda$ controls the orthogonal constraint. While computationally efficient, these models occasionally fail to capture the high-frequency volatility jumps inherent to tick-level financial data.

\subsection{\revadddec{Imputation: }Adversarial and Generative Imputation Dynamics}

To overcome the smoothing effects typical of standard interpolators, Generative \revdeldec{ Adversarial Networks }\revadddec{adversarial networks }(GANs)\revadddec{~\citebox{\citep{goodfellow2014gan,yoon2018gain}} }formulate missing data recovery as a competitive game. However, standard GANs frequently suffer from severe distributional incoherence—commonly referred to as "drift"—during contiguous data outages. To directly address this, AttnWGAIN \citep{wang2025attn} substitutes traditional convolutional layers with diagonal masked attention blocks, explicitly forcing the generator to jointly learn temporal and feature-wise correlations without accessing masked future tokens. Similarly, ImputeGAN~\citep{qin2023imputegan} leverages an Informer-based generator\revadddec{~\citep{zhou2021informer} }utilizing probabilistic sparse attention, significantly reducing the $O(L^2)$ computational complexity of long-sequence dependencies while maintaining fidelity.

To rigorously enforce distributional stability across extensive missing windows, the CWGAIN-GP framework \citep{wang2024cwgain} employs a \revdeldec{ Conditional }\revadddec{conditional }Wasserstein GAN fortified with a \revdeldec{ Gradient Penalty}\revadddec{gradient penalty~\citebox{\citep{arjovsky2017wgan,gulrajani2017wgangp}}}. By conditioning the adversarial game on the observation mask $\mathbf{M}$ and enforcing Lipschitz continuity, it successfully bounds the generative drift, ensuring that the synthetic interpolations mathematically align with the heavy-tailed distributions of surrounding observable financial regimes.

\subsection{\revadddec{Imputation: }Diffusion and Foundation Models for Missing Data}

Score-based diffusion models\revdeldec{ have recently established a new state-of-the-art in financial }\revadddec{~\citebox{\citep{ho2020ddpm,song2019generative,song2021score}} have become a leading approach for probabilistic }imputation by treating missing value reconstruction as a conditional generative process. These architectures define a forward process that adds Gaussian noise strictly to the missing regions (or the entire sequence, masked subsequently), and optimize a reverse denoising step conditioned on $\mathbf{X}_{obs}$:
\begin{equation}
\mathcal{L}_{diff} = \mathbb{E}_{\mathbf{x}_0, \boldsymbol{\epsilon}, t}[\|\boldsymbol{\epsilon} - \boldsymbol{\epsilon}_\theta(\mathbf{x}_t, t, \mathbf{X}_{obs})\|^2]
\end{equation}
A \revdeldec{ seminal }\revadddec{central }architecture in this domain is CSDI \citep{tashiro2021csdi}, which utilizes a dual 2D-attention mechanism to seamlessly integrate temporal dynamics and cross-feature correlations. Subsequent works such as LSSDM, Score-CDM, and SaSDim \citep{LSSDM, Score-CDM, SaSDim} further advanced this paradigm. To eliminate the "boundary disharmony" problem—unrealistic volatility jumps at the borders of observed and imputed regions—advanced diffusion models incorporate explicit consistency constraints across temporal and spatial horizons \citep{zhang2024spatial}, as well as cross-variable boundaries \citep{zhou2024cross}.

Furthermore, models like LSCD \citep{fons2025lscd} and SPDM \citep{SPDM} employ frequency-domain transformations to map 1D time series into 2D periodic representations. This isolates stable macroeconomic cyclical components from high-frequency market microstructure noise during the denoising phase. Expanding this to transformer scales, TimeDiT \citep{cao2024timedit} introduces a \revdeldec{ Diffusion }\revadddec{diffusion }Transformer (DiT)\revadddec{~\citep{peebles2023scalable} }utilizing a unified masking strategy, uniquely enabling fine-tuning-free model editing. This allows quantitative analysts to directly inject external physical constraints or market heuristics into the sampling trajectory.

Simultaneously, \revdeldec{ Large Language Models }\revadddec{large language models }(LLMs)\revdeldec{ have been ingeniously }\revadddec{, building on few-shot language-model capabilities~\citep{brown2020language}, have been }adapted for heterogeneous tabular and sequential data imputation. LLM-Forest \citep{chen2023llmforest} pioneers an ensemble approach utilizing graph-augmented prompts. It constructs bipartite graphs to map data entries to feature values, utilizing hierarchical merging to retrieve nearest neighbors. These neighbors serve as few-shot prompts for a forest of LLMs, resolving the final imputed value via confidence-weighted voting. Additionally, frameworks like MTabGen \citep{liu2023mtabgen} bridge diffusion dynamics and mixed-type tabular formats by running parallel forward processes for categorical and continuous numerical features, unifying them via columnar embeddings.

\subsection{\revadddec{Imputation: }Emerging Trends}

Recent literature \revdeldec{ underscores a definitive paradigm }\revadddec{indicates a substantial }shift toward handling complex, non-random missing patterns through specialized extensions of these generative frameworks. As \revdeldec{ comprehensively }documented in recent methodological surveys by \revadddec{Shan et al.~}\citep{shan2021} and \revadddec{Du et al.~}\citep{du2024}, the field is rapidly moving away from isolated point-interpolators toward unified, context-aware sequence modelers. We highlight three critical emerging trends driving this evolution.

The first trend concerns unified multi-task architectures: instead of treating imputation as an isolated preprocessing step, modern models increasingly solve related financial tasks simultaneously to improve representational learning. For example, \revadddec{Miao et al.~}\citep{miao2025} propose coupled diffusion models that jointly perform forecasting and imputation, leveraging shared latent representations to enforce temporal consistency. To enhance structural robustness, \revadddec{Zhang et al.~}\citep{zhang2023imputation} \revdeldec{ utilize }\revadddec{use }dual-attention networks to simultaneously capture local and global dependencies, while \revadddec{Cai et al.~}\citep{cai2025} employ \revdeldec{ Temporal Graph Networks }\revadddec{temporal graph networks }(TGNs)\revadddec{~\citep{rossi2020temporal} }to map dynamic cross-asset spillovers during fragmented market outages \citep{chen2023b}. Furthermore, \revadddec{Zhang et al.~}\citep{zhang2025a} extend this paradigm beyond numerical data, advancing multi-modal frameworks that integrate textual sentiment to inform the imputation of missing price data.

The second trend focuses on high-frequency and microstructure imputation: it centers on the granular complexities of high-frequency trading (HFT) and Limit Order Book (LOB) data. Standard daily-level imputers often fail at the microstructural level; thus, \revadddec{Bamford et al.~}\citep{bamford2023} and \revadddec{Shankar et al.~}\citep{shankar2025} have proposed specialized deep learning architectures and benchmark suites that strictly respect LOB queue dynamics and volume imbalances. To bridge the gap between generative fidelity and the stringent latency requirements of HFT systems, \revadddec{Nater and Khayati~}\citep{nater2025} introduce scalable probabilistic imputation techniques that bypass the computational overhead of iterative diffusion without sacrificing distributional shape.

The third trend concerns regime-awareness and stress-testing protocols: the evaluation of imputation models has matured to account for market non-stationarity. Recognizing that structural breaks often coincide with missing data (e.g., trading halts), \revadddec{Qian et al.~}\citep{qian2024} and \revadddec{Gao et al.~}\citep{gao2024} introduce dynamic, regime-aware score-based models specifically designed to adapt to shifting market volatility. Crucially, a growing body of work advocates for \textit{stratified stress-testing evaluation protocols}. \revadddec{Wi et al.~}\citep{wi2023} emphasize that imputation fidelity must be validated during simulated market crashes rather than strictly during stationary periods, ensuring that imputed data maintains structural integrity when risk managers need it most. Complementing this \revdeldec{ rigorous evaluation, }\revadddec{evaluation, Wang et al.~}\citep{wang2024c} \revdeldec{ formally }investigate the pervasive issue of look-ahead bias in deep learning imputers, establishing strict cross-validation guidelines to ensure that reconstructed historical sequences do not inadvertently leak future market information into backtesting pipelines.

\subsection{\revdeldec{ Summarization}\revadddec{Imputation: Research Synthesis}}

The landscape of financial missing data recovery is defined by a rigorous trade-off between computational efficiency and generative fidelity, as synthesized in Table~\ref{tab:imputation_comparison}. The empirical case for diffusion-based imputation is, however, less settled than the tabular taxonomy above might suggest. The most-cited diffusion imputer, CSDI~\citep{tashiro2021csdi}, reports a 40--65\% improvement over probabilistic baselines and a 5--20\% improvement over deterministic baselines---but on \emph{healthcare and environmental} datasets, not financial. Direct head-to-head benchmarks that put diffusion imputers (CSDI, TimeDiT, LSCD) on the same financial dataset as GAN imputers (GAIN, AttnWGAIN, ImputeGAN, CWGAIN-GP) are not standard in the literature, and the few within-paper comparisons that exist (e.g.\ Score-CDM vs GAIN on air-quality or electricity data) do not transfer cleanly to the financial regime of contiguous missingness and regime-switching volatility. The honest reading is therefore: diffusion models are the most promising current direction for high-fidelity imputation under contiguous missingness, but the claim that they ``stand unrivaled'' for financial imputation specifically rests on architectural plausibility and on transferred evidence rather than on a published financial benchmark. Tensor factorization and tree-based methods remain preferable when the constraint is large-scale historical-database throughput rather than peak reconstruction quality.
